\chapter{摩尔斯电码}
% 完整版：chapters/04_code.tex

摩尔斯电码有两个符号——点和划——以及它们之间的停顿。神经元的字母表更加贫乏：只有一个符号，短促的咔嗒。所以，整条消息都写在停顿里。神经元之间说的究竟是什么语言？

\section*{咔嗒能说多少话}

先像通信工程师那样发问：这样的信道到底能装下多少信息？如果接收者能以毫秒精度分辨咔嗒到达的时刻，一次咔嗒大约能携带八比特。如果接收者只数一秒内来了几次咔嗒，那么从同一串信号里，它提取的信息要少几十倍。信道容量几乎全都藏在\emph{时间}里。

\pencilfig{4}{\pencilart{4}}{上面是摩尔斯电码，下面是神经元：只有一个符号，咔嗒；整条消息都写在停顿里。}

但信道有噪声。说来奇怪，神经元本身是精确的：把同一段变化的电流反复输入给它，它每次都会在相同的时刻咔嗒，误差不超过一毫秒。噪声来自连接：跑到导线末端的咔嗒，有一半以上直接丢失了——并不是每次都会释放物质。在活的皮层里，咔嗒流看上去几乎是随机的。这是噪声，还是我们还读不懂的消息？压缩得好的消息，在局外人看来和随机数没有区别，所以这个问题仍然悬而未决。

\section*{慢，但可靠}

最简单的编码是数咔嗒：越频繁，数越大。这种编码不怕丢失，也不怕抖动。但它慢。要把频率测到百分之十的精度，需要一百次咔嗒，也就是好几秒。可是，一张图片在你眼前一闪而过，你大约只用十分之一点五秒就能认出上面的动物，而信号在此期间要经过大约十级处理。在每一级上，每个神经元至多来得及咔嗒一次。

出路是同时用许多神经元：不是一个听众数一分钟，而是一百个听众数一瞬间。可这里也有陷阱：相邻神经元的噪声有些相似，就像人群里的传言。对几十个邻居取平均是有用的，再往上就几乎没用了。

\pencilfig{122}{%
% error of the group-average rate relative to one neuron, sqrt(c + (1-c)/N): c = 0.1 (neighbours' noise
% is slightly alike; full edition, chapter 4) and c = 0 (independent noise). x = 0.8 log2 N, y = 2.2 * error
\draw[pen] (0,0) -- (5.9,0); \draw[pen] (0,0) -- (0,2.45);
\foreach \x/\n in {0/1, 2.66/10, 5.31/100}{\draw[pcGray, line width=0.5pt] (\x,0) -- (\x,-0.08); \node[lbl, anchor=north] at (\x,-0.08) {\n};}
\node[lbl, anchor=north] at (2.95,-0.38) {参与平均的神经元数};
\node[lbl, anchor=south, rotate=90] at (-0.1,1.2) {误差};
\draw[pcGray, densely dashed, line width=0.5pt] (0,0.70) -- (5.6,0.70);
\draw[penlite=pcBlue] plot[smooth] coordinates {(0.00,2.20) (0.47,1.80) (0.80,1.56) (1.27,1.27) (1.60,1.10) (2.07,0.90) (2.40,0.78) (2.87,0.64) (3.20,0.55) (3.67,0.45) (4.00,0.39) (4.47,0.32) (4.80,0.28) (5.27,0.22) (5.60,0.19)};
\draw[penlite=pcRed] plot[smooth] coordinates {(0.00,2.20) (0.47,1.84) (0.80,1.63) (1.27,1.39) (1.60,1.25) (2.07,1.10) (2.40,1.01) (2.87,0.92) (3.20,0.87) (3.67,0.82) (4.00,0.79) (4.47,0.76) (4.80,0.74) (5.27,0.73) (5.60,0.72)};
\node[lbl, text=pcRed, anchor=west, align=left] at (5.7,0.74) {噪声相似，\\如皮层};
\node[lbl, text=pcBlue, anchor=west] at (5.7,0.19) {噪声独立};
\node[lbl, anchor=west] at (0.9,2.15) {单个神经元};
\node[lbl, anchor=north west] at (0.05,0.66) {只剩三分之一};
}{对一群神经元取平均后的频率误差。如果噪声相互独立，误差会一直下降。邻居的噪声相似，误差就停在大约三分之一处，而且几十个神经元就已经到了这个极限。}

\section*{快：谁先到}

第二条出路是把数写进时间里。强信号让神经元早早咔嗒，弱信号让它晚些咔嗒。仅仅一次咔嗒就已经携带了一个数。可是，没有公共的起点，时间从哪里算起？可以让神经元彼此比较：谁第一个咔嗒，谁第二个。十个神经元咔嗒到达的先后顺序就携带二十多比特。此外，一次齐射也可以充当起点标记——就像摩尔斯电码里单词之间的长停顿。

\pencilfig{121}{%
% latency code: one click per neuron; the stronger the input, the earlier the click
\foreach \y/\t/\l/\k in {1.5/1.1/{强输入}/1, 0.95/2.4/{中等}/2, 0.4/4.2/{弱}/3}{
  \draw[pen] (0.4,\y) -- (5.1,\y);
  \draw[pcBlue, line width=0.9pt] (\t,\y) -- (\t,\y+0.42);
  \node[lbl, anchor=east] at (0.3,\y+0.1) {\l};
  \draw[dotpen=pcAmber, wash=pcAmber] (5.6,\y+0.16) circle (0.17);
  \node[font=\scriptsize, text=pcAmber!60!black] at (5.6,\y+0.16) {\k};}
\draw[pcGray, densely dashed, line width=0.5pt] (0.55,0.25) -- (0.55,2.1);
\node[lbl, anchor=south] at (0.55,2.1) {刺激};
\node[lbl, text=pcAmber!70!black, anchor=south] at (5.6,2.1) {顺序};
\draw[arr] (0.7,0.05) -- (4.9,0.05); \node[lbl, anchor=north] at (2.8,0.02) {时间};
}{强输入——咔嗒早，弱输入——咔嗒晚。即使不知道起始时刻，咔嗒的顺序也能说明哪个输入更强。}

\section*{编码由听者选择}

同一串咔嗒既携带频率，也携带时刻和顺序。其中哪一样会被读出来，由接收者决定。桶上“洞很宽”的神经元慢慢攒水，只听得见频率。遗忘得快的神经元只听得见同时到达。而刹车，正如我们所见，能把神经元从一种模式切换到另一种。

就连连接本身也会过滤信号。有些连接很快就疲劳，因此报告的不是频率，而是频率的\emph{变化}。另一些则相反，更容易让连续几次咔嗒组成的簇放电通过：神经元于是有了自己的“划”。与此同时，抑制性神经元负责打标点——把信号流切成一个个窗口，并在它太响时把音量压低。

\section*{节俭的语言}

最后一点：这种语言很昂贵。每一次咔嗒都要消耗能量，而整个大脑的功耗只有二十瓦左右。把这些瓦数平摊到所有神经元上，平均每个神经元每秒大约只能咔嗒一次。大多数神经元在大部分时间里是沉默的。大脑的编码必须吝啬。

在摩尔斯电码里，英语最常见的字母 E 就是一个点：常见的东西用短码。神经元也差不多：意料之中的事只花很少的咔嗒。持续不变的刺激渐渐不再引起响应，传感器报告的是变化，而多巴胺只报告意外。大脑把咔嗒花在新闻上。

\fullversion{4}
